\documentclass{umk-matematika}

\begin{document}

\maketitlepage
    {Практическая работа №14}
    {Двугранный угол}
    {}

\section{Фундамент (1--4 балла)}

\textbf{Задача 1 (1 балл).} Что такое двугранный угол?

\textbf{Задача 2 (2 балла).} Что такое линейный угол двугранного угла?

\textbf{Задача 3 (3 балла).} Чему равен двугранный угол между стеной и полом в комнате?

\textbf{Задача 4 (4 балла).} Какие плоскости называются перпендикулярными?

\section{Стены (5--7 баллов)}

\textbf{Задача 5 (5 баллов).} В кубе $ABCDA_1B_1C_1D_1$ найдите линейный угол двугранного угла между плоскостями ABC и AB₁C₁.

\textbf{Задача 6 (6 баллов).} Двугранный угол равен 60$^\circ$. На одной грани взята точка A, удалённая от ребра на 10 см. Найдите расстояние от точки A до другой грани.

\textbf{Задача 7 (7 баллов).} Признак перпендикулярности двух плоскостей.

\section{Кровля (8--10 баллов)}

\textbf{Задача 8 (8 баллов).} В правильной треугольной призме сторона основания равна 4 см, боковое ребро — 6 см. Найдите двугранный угол между плоскостью основания и плоскостью, проходящей через сторону основания и середину противолежащего бокового ребра.

\textbf{Задача 9 (9 баллов).} Двугранный угол равен 120$^\circ$. Точка A лежит внутри угла и удалена от обеих граней на расстояние 4 см. Найдите расстояние от точки A до ребра двугранного угла.

\textbf{Задача 10 (10 баллов).} Крыша имеет форму двугранного угла 120$^\circ$. На скате крыши на расстоянии 3 м от конька (ребра) установлена опора, перпендикулярная скату. На какой высоте от основания крыши находится точка крепления опоры, если ширина дома 8 м?

\newpage
\section*{Ответы}

\begin{enumerate}
  \item две полуплоскости с общей границей.
  \item угол между перпендикулярами к ребру.
  \item 90$^\circ$.
  \item двугранный угол = 90$^\circ$.
  \item 90$^\circ$.
  \item $5\sqrt{3}$ см.
  \item признак перпендикулярности плоскостей.
  \item $60^\circ$.
  \item 8 см.
  \item ≈ 4,9 м.
\end{enumerate}

\end{document}